\section{Introduction}
\label{sec:introduction}

% ---------------------------------------------------------------
%  PARAGRAPH 1: MOTIVATION — Why QSP/QSVT matters
% ---------------------------------------------------------------

Quantum signal processing
(QSP)~\cite{LowChuang2017,GSLW2019} and its extension to quantum
singular value transformation
(QSVT)~\cite{GSLW2019} have emerged as a unifying framework for
quantum algorithms.  By interleaving a signal unitary with tunable
phase rotations, QSP constructs polynomial transformations of a
quantum operator's eigenvalues, subsuming Hamiltonian
simulation~\cite{LowChuang2017,BerryChildsCleveEtAl2015}, quantum
linear-systems solvers~\cite{HHL2009}, eigenvalue filtering, and a
wide range of linear-algebraic primitives under a single algorithmic
paradigm~\cite{MartynRossiTanChuang2021}.  The framework's power
stems from a precise characterization of the achievable function
class: every bounded polynomial of definite parity can be realized,
and the corresponding phase angles can be found
efficiently~\cite{DongMengWhaleyEtAl2021,Haah2019}.

% ---------------------------------------------------------------
%  PARAGRAPH 2: THE BOUNDEDNESS CONSTRAINT — What it costs
% ---------------------------------------------------------------

A fundamental limitation of this framework, however, is
\emph{boundedness}.  The signal value $a = \cos(\theta/2)$ lies in
$[-1,1]$, and the output polynomial must satisfy $|P(a)| \leq 1$ for
all $a$ in that interval---a direct consequence of
unitarity~\cite{LowChuang2017}.  When QSP is applied via block
encoding to a Hermitian operator~$H$, this constraint forces a
normalization $H \to H/\alpha$ with $\alpha \geq \|H\|$, which
requires knowledge of the operator norm, increases circuit depth by a
factor of~$\alpha$~\cite{GSLW2019,MartynRossiTanChuang2021}, and
restricts the target transformation to functions approximable by
polynomials bounded by~$1$ on $[-1,1]$---excluding unbounded functions
such as $1/\lambda$, $\lambda$, or $\tan(\lambda)$ without truncation
and regularization.  That boundedness is a genuine restriction, not merely
a technical inconvenience, was formalized by Tang and
Tian~\cite{TangTian2024}, who exhibit functions well-approximable by
unbounded polynomials on $[-1,1]$ but provably inapproximable by
bounded ones.

% ---------------------------------------------------------------
%  PARAGRAPH 3: LANDSCAPE — Recent QSP advances (none lift boundedness)
% ---------------------------------------------------------------

Recent work has extended the QSP framework in several directions.  One
line of development broadens the achievable function class: generalized
QSP (GQSP)~\cite{MotlaghWiebe2023} replaces $Z$-rotation signal
processing operators with full $\SU(2)$ rotations, lifting the parity
and realness constraints, while Quantum Phase
Processing~\cite{WangZhangYuWang2023} applies trigonometric
transformations to eigenphases and Rossi, Ceroni, and
Chuang~\cite{RossiCeroniChuang2025} extend QSP to composable
multivariable polynomial transforms.  A second line reduces the
classical or quantum cost of the framework:
Alase~\cite{Alase2025} bypasses the angle-finding bottleneck by
constructing QSP circuits directly from function evaluations,
Randomized QSVT~\cite{WangZhangHazraEtAl2025} avoids block encodings
through importance sampling, and Kang and Su~\cite{KangSu2025}
introduce Hamiltonian block encoding using at most two ancilla qubits
without explicit normalization.  These advances enrich the QSP toolkit
in expressiveness, classical preprocessing cost, and circuit depth,
yet all retain a common limitation: the output polynomial
remains bounded, $|P| \leq 1$, and reads out a single amplitude
rather than a ratio of amplitudes.

% ---------------------------------------------------------------
%  PARAGRAPH 4: APPROACH — What this paper does
% ---------------------------------------------------------------

In this work we ask whether the boundedness constraint on the
\emph{decoded output} can be circumvented without modifying the
quantum algorithm itself.  We show that it can, through a change
of the classical encoding and decoding interface.
The key idea is to compose QSP with \emph{stereographic encoding}---the
conformal bijection between the extended complex plane
$\hat{\C} = \C \cup \{\infty\}$ and the Bloch
sphere~$\Sphere$.  On the input side, an unbounded parameter
$r \in [0,\infty)$ is compressed into the bounded QSP domain via
$\rtilde = r/\sqrt{1+r^2} \in [0,1)$, so that QSP operates
internally on a bounded signal.  On the output side, decoding the QSP
output state via Pauli measurements produces the \emph{ratio}
$P(\rtilde)/Q(\rtilde)$ of the two state components---a rational
function that diverges wherever $Q(\rtilde) = 0$.  The external
input--output relationship $r \mapsto f(r)$ is therefore an unbounded
rational function of the original unbounded parameter, even though QSP
operates internally on bounded values.  The unitarity constraint
$|P(\rtilde)|^2 + |Q(\rtilde)|^2 = 1$ is preserved, but it no longer
implies boundedness of the output: the ratio of two bounded quantities
need not be bounded.  A key structural result is that stereographic
and standard QSP signal operators are related by the basis change
$V = \mathrm{diag}(1,i)$, so the entire machinery of standard
QSP---phase-finding algorithms, function-class characterization,
convergence theory---applies directly, with no new algorithms required.

% ---------------------------------------------------------------
%  PARAGRAPH 5: SUMMARY OF RESULTS
% ---------------------------------------------------------------

Our main contributions are as follows.  We define the stereographic
encoding $\ket{z} = (z\ket{0}+\ket{1})/\sqrt{|z|^2+1}$ and show
that single-qubit gates act as M\"obius transformations on the encoded
value (Section~\ref{sec:encoding}).  We then prove that the $k$-fold
stereographic signal operator produces the state
$T_k(\rtilde)\ket{0} + U_{k-1}(\rtilde)/\!\sqrt{1+r^2}\,\ket{1}$,
where $T_k$ and $U_{k-1}$ are Chebyshev polynomials whose compositions
with the stereographic map are Boyd's rational Chebyshev
functions~\cite{Boyd1990}---an orthogonal basis on
$(-\infty,\infty)$ with weight $1/(1+r^2)$.  Decoding this state
yields the unbounded rational function
$z_k(r) = \cot(k\arctan(1/r))$, with $k$ interlacing zeros and poles
on $[0,\infty)$.  The stereographic and standard QSP signal operators
are unitarily equivalent via $V = \mathrm{diag}(1,i)$, immediately
establishing the full function class and the existence of phase
factors from the standard QSP theorem (Section~\ref{sec:main}).
We quantify the cost tradeoff: the coherent normalization overhead of
standard QSP is replaced by a polynomial sampling cost, with
absolute precision $\epsilon$ requiring
$\mathcal{O}(r^6/\epsilon^2)$ shots
(Section~\ref{sec:cost-analysis}).  The extension to multi-qubit
eigenvalue transformations via stereographic block encoding is
developed in a companion paper~\cite{KharaziLeytonOrtega2026b}.

% ---------------------------------------------------------------
%  PARAGRAPH 6: PAPER ORGANIZATION
% ---------------------------------------------------------------

The remainder of this paper is organized as follows.
Section~\ref{sec:preliminaries} reviews the standard QSP framework,
the role of Chebyshev polynomials, and the boundedness constraint.
Section~\ref{sec:encoding} introduces stereographic encoding as a
geometric alternative to amplitude encoding, deriving its decoding
formula and gate--M\"obius correspondence.
Section~\ref{sec:main} contains the main results: the stereographic
signal operator, Chebyshev state generation, the unbounded rational
function class, and the basis-change reduction to standard QSP.
Section~\ref{sec:cost-analysis} quantifies the depth-versus-sampling
tradeoff, compares with standard and generalized QSP, and presents
numerical results for four target functions.
Section~\ref{sec:applications} develops an application to quantum
state preparation, showing how stereographic QSP extends recent
QSVT-based methods~\cite{McArdleGilyenBerta2025} to unbounded and
rational target functions.
Section~\ref{sec:discussion} discusses connections to prior work
and open problems.  The overall framework is illustrated
schematically in Fig.~\ref{fig:schematic}.
